\section{Structural alternatives, measured}
\label{sec:alternatives}

Every row above is engineering against a fixed exponent. Three
threads asked whether the exponent could move for this curve, in the
same accounting, and each is reported here with its measurement
because a later reader who rediscovers the idea should find the
number rather than the folklore. None changed the campaign.

\subsection{Index calculus on the Koblitz curve}
\label{sec:alt-ic}

Index calculus on $E(\FF_{2^n})$ decomposes a random point into $m$
factor-base points through Semaev's summation
polynomials~\cite{semaev2004,gaudry2009,diem2011,faugere2012,petitquisquater2012,galbraithgebregiyorgis2014};
on a Koblitz curve the factor base can be taken Frobenius-stable so
that columns are orbits, saving a factor $n$ in relations and $n^2$ in
linear algebra against rho's $\sqrt{2n}$. Two structural facts about
$n=131$ come first. In a normal basis the Frobenius is a coordinate
permutation, so a Frobenius-stable subspace is an
$\FF_2[t]/(t^n-1)$-submodule, a binary cyclic code of length $n$, and
at $n=131$ (and $n=163$) $2$ is primitive modulo $n$, so $t^n-1$ has
exactly two factors and the only stable subspaces have dimension
$0,1,130,131$. The usual subspace construction is empty at exactly the interesting
sizes; forced onto the 130-dimensional subspace, the orbit trick is a
net loss of $57$--$96$ bits for
$m\in\{2,\ldots,5\}$.\art{RESEARCH_KOBLITZ_INDEX_CALCULUS.md}\art{autoresearcher:ledger/evidence/EV-FROB-d336b0.yaml}
The Couveignes--Lercier elliptic-period alternative needs an auxiliary
curve over $\FF_2$ with $n\mid\#E'(\FF_2)$, and Hasse gives at most
5.\art{RESEARCH_KOBLITZ_INDEX_CALCULUS.md} What survives is a
weight-bounded base $\{R:\HW(x_R)\le w\}$ in the normal basis, which
is Frobenius-stable for every $n$, is the predicate the walk already
computes, and is not a low-degree algebraic condition, so only a
solver that does not need linearity can use it.

\paragraph{SAT decomposition.} Chaining Semaev's $S_3$ through the
intermediate sums, with squaring free, encodes a $k$-point
decomposition as a CNF whose multiplications cost the generated
Karatsuba circuit's gates ($15{,}588$ at $n=131$ against $2.2$ million for the definition). With $k=3$, $w=3$ and CryptoMiniSat the median solve is $0.11$~s at
$n=9$ and $2.64$~s at $n=11$; at $n=23$ no instance finished, one
running 25 minutes without returning: roughly $2^{2.3}$ per field bit
over that range, from three points.\art{ecc2k130/README.md, ``Index calculus: SAT point decomposition''}

\paragraph{The compact-orbit ladder.} A separate line of work solves rungs of the same curve family ($a=0$,
the ECC2K-130 equation) with a four-summand decomposition, at $n=41$
and $53$ over 85- and 220-column bases and at $n=61,71,73$ over a
600-column base, fully charged, and fits the rank cost. Table~\ref{tab:ic-rungs}
is the measured cost and its extrapolation to $n=131$ under the fit
$\text{probes/relation}\approx C\,r/(n^2K^2)$ with
$C=1$.\art{RESEARCH_ECC2K130_IC_FEASIBILITY.md} The fitted $C$ ($0.33$ at $n=61$, $1.51$ at $n=71$, $1.27$ at $n=73$,
recomputed here from the exact subgroup orders) disagrees with the
structural model's $0.25$ by a factor of $1.3$ to $6$, and the note
records the gap as open; the extrapolation uses $C=1$, between the
structural $0.25$ and the fitted $1.3$--$1.5$ of the two largest
rungs.

\begin{table}[t]
\centering
\small
\caption{Four-summand compact-orbit index calculus on the ECC2K-130
  curve family, single thread, $K=600$ columns, precompute fully
  charged (left), and the extrapolation to $n=131$ ($r=6.8\cdot10^{38}$)
  at $1.85\cdot10^6$ probes/s (right). The minimum base for one
  expected four-tuple per point is $B>r^{1/4}=5.1\cdot10^9$ points,
  $K\ge2\cdot10^7$.}
\label{tab:ic-rungs}
\begin{tabular}{@{}rrrr@{}}
\toprule
$n$ & $r$ (bits) & rank wall & probes/relation \\
\midrule
61 & 47.2 & $13.1$~s & $\sim4\cdot10^4$ \\
71 & 52.3 & $1{,}485$~s & $\sim4.6\cdot10^6$ \\
73 & 56.3 & $18{,}398$~s & $5.7\cdot10^7$ \\
\bottomrule
\end{tabular}
\hspace{1.5em}
\begin{tabular}{@{}rrrr@{}}
\toprule
$K$ & probes/relation & total probes & core-years \\
\midrule
$600$ & $1.1\cdot10^{29}$ & $6.6\cdot10^{31}$ & $1.1\cdot10^{18}$ \\
$10^6$ & $4.0\cdot10^{22}$ & $4.0\cdot10^{28}$ & $6.8\cdot10^{14}$ \\
$10^9$ & $4.0\cdot10^{16}$ & $4.0\cdot10^{25}$ & $6.8\cdot10^{11}$ \\
$10^{12}$ & $4.0\cdot10^{10}$ & $4.0\cdot10^{22}$ & $6.8\cdot10^{8}$ \\
\bottomrule
\end{tabular}
\end{table}

The verdict is in the unit of the rest of the paper. Rho at $2^{60.9}$ iterations on one RTX~PRO~6000 is about five
GPU-years at the campaign's $14.11$~B/s. Four-summand index calculus
has total work about $r/(n^2K)$ probes: at $K=10^9$ that is
$4\cdot10^{25}$, some $2\cdot10^7$ times rho's iteration count with one
probe counted as one step (uncalibrated), and the base at which the precompute alone equals
rho is $K\approx1.8\cdot10^{16}$ columns, beyond any linear algebra
contemplated. The rungs that beat rho do so only in the online class
after a precompute that is excluded from the clock (at $n=73$, $18{,}400$~s of precompute against $46$~s of rho online, a
$400\times$ total-work deficit traded for a $185\times$ exploratory
online wall ratio on one frozen target, on a shared host), and with
the precompute charged the method spends $26$--$1{,}200\times$ more
probes than one rho solve takes steps on every rung, again
uncalibrated.
Beating rho on total work needs a relation-shape change (five- or
six-summand relations, whose per-point yield scales as $B^5/(120r)$),
not engineering, and that oracle is unmeasured.\art{RESEARCH_ECC2K130_IC_FEASIBILITY.md, §3--4}
In this repository's boundary ledger the Koblitz rows at a 39-bit
subgroup run from $S=\text{ops}/\sqrt r=58.8$ ($300\times$
signed-Frobenius rho) for the plain three-summand meet-in-the-middle
oracle down to $5.12$ ($26\times$; $5.92$ and $30\times$ on
re-measurement) for a balanced two-summand base; the secondary
complete-CPU panel on disjoint public points at $n=37,41,53$ has rho
ahead in every cell.\art{research/notes/index-calculus/RESEARCH_IC_BOUNDARY_LEDGER.md}\art{docs/ic/BOUNDARY_TARGETS.md}

\subsection{Point decomposition against the per-call budget}
\label{sec:alt-pdp}

The sibling library prices the same question from the other end:
given rho's $2^{60.81}$ orbit-walk iterations, how fast must one
$m$-summand decomposition be for index calculus to fit inside that
budget, and how fast is the best one measured? With a base of
$2^\ell$ points, $m=4$ at $\ell=29$ needs $2^{48.6}$ decompositions and
$2^{60.0}$ of linear algebra, leaving $2^{12}$ rho-iterations per
decomposition; $m=5$ at $\ell=28$ leaves $2^{33}$; $m=6$ at $\ell=24$
leaves $2^{37}$; $m=3$ has no budget at all because the linear algebra
alone exceeds rho. The best measured decomposition is the
non-algebraic meet-in-the-middle at exponent $\lceil m/2\rceil$ in
$\ell$, and extrapolated to $n=131$ it is $2^{58}$, $2^{84}$ and $2^{72}$
for $m=4,5,6$: $35$--$51$ bits over the budget. Algebraic solvers land far higher: SAT fits at $2^{4}$ ($m=3$) to
$2^{7}$ ($m=4$) per unit of $\ell$, msolve grows faster still, and
WDSat at about $2^{3}$, the exhaustive-search exponent. An attack needs a decomposition about $2^{35}$--$2^{50}$
times faster than meet-in-the-middle at $\ell\approx25$--$30$, which is
the obstruction that stops Wagner's $k$-tree from applying to elliptic
curves; the trace-homomorphism clause that halves UNSAT conflicts is a
constant, not an exponent.\art{cryptanalysis:experiments/pdp-scaling/README.md}
A Riemann--Roch (Nagao) oracle on the real curve was measured at
factor-base dimensions 4--8 and its derived product law
($\text{relations}\times\text{targets}\times\text{oracle}=m\cdot2^{131}$)
leaves even a free oracle above rho at $m=3$ ($2^{68.6}$) and below it
only from $m=4$ ($2^{56.4}$, leaving $2^{4.4}$ of headroom for the
oracle itself), where no measured oracle comes close.\art{autoresearcher:ledger/evidence/EV-ICPERF-10c5fc.yaml}

\subsection{What the negatives say}

Each alternative was carried to a measurement and each landed above
rho by a margin that is a change of exponent, not of constant. The
honest summary for ECC2K-130 is that the generic-group floor,
$\sqrt{\pi\ell/(2\cdot262)}$ class-walk steps, is still the number
the campaign is walking toward, and that the only lever that has
moved in this work is the cost of one step.
